% ==================== 第 4 章 总体设计 ====================
\section{总体设计}

\subsection{设计原则}

本项目在系统设计时遵循以下原则：
\begin{enumerate}[leftmargin=2em,itemsep=0.2em]
\item \textbf{算法与界面解耦：}所有业务数据通过 \code{io.py} 转换为
纯 Python 对象（\code{Problem}），由 \code{solver.py} 求解，
与具体 UI 实现无关。
\item \textbf{硬约束求解与软目标拆解：}硬约束用 \code{AddAtMostOne} /
\code{AddExactlyOne} 保证，软目标用线性加权和实现，
便于对比与解释。
\item \textbf{零依赖部署：}仅依赖 OR-Tools 单一第三方库，
HTTP 服务使用 Python 标准库，前端无构建步骤。
\item \textbf{可测试性：}每个模块可独立 \code{unittest} 覆盖，
业务逻辑通过纯函数实现，避免副作用。
\end{enumerate}

\subsection{系统架构}

\begin{figure}[h]
\centering
\begin{tikzpicture}[
  node distance=0.7cm and 0.9cm,
  layer/.style={font=\bfseries\color{cmblue}},
  every node/.style={font=\small}
]
\node[layer] (fe) {前端层 (Browser)};
\node[cmblock, below=of fe] (dash) {决策驾驶舱};
\node[cmblock, right=of dash] (sched) {智能排课};
\node[cmblock, below=of dash] (resc) {调课中心};
\node[cmblock, right=of resc] (reso) {资源概览};

\node[layer, below=1.2cm of dash.south] (be) {后端层 (Python 3.12)};
\node[cmblock, fill=cmgreen!10, below=0.7cm of be] (api) {API 网关 (api.py)};
\node[cmblock, fill=cmgreen!10, below=of api] (svc) {业务编排 (service.py)};

\node[layer, below=1.2cm of svc.south] (core) {核心算法层};
\node[cmblock, fill=cmaccent!10, below=0.7cm of core] (solv) {求解器 (solver.py)};
\node[cmblock, fill=cmaccent!10, right=of solv] (vali) {校验器 (validator.py)};
\node[cmblock, fill=cmaccent!10, left=of solv] (io) {数据 I/O (io.py)};

\node[cmblock, below=of vali, fill=gray!10] (mdl) {models.py (数据类)};

\node[layer, below=0.7cm of mdl.south] (data) {数据层};
\node[cmblock, below=0.7cm of data] (demo) {data/demo.json};

\draw[cmarrow] (dash) -- (api);
\draw[cmarrow] (sched) -- (api);
\draw[cmarrow] (resc) -- (api);
\draw[cmarrow] (reso) -- (api);
\draw[cmarrow] (api) -- (svc);
\draw[cmarrow] (svc) -- (solv);
\draw[cmarrow] (svc) -- (vali);
\draw[cmarrow] (solv) -- (io);
\draw[cmarrow] (io) -- (demo);
\draw[cmarrow, dashed] (vali) -- (mdl);
\draw[cmarrow, dashed] (solv) -- (mdl);
\end{tikzpicture}
\caption{ClassMind 系统架构图}
\label{fig:arch}
\end{figure}

\subsubsection{前端层}

前端由 4 个独立子页面 + 1 个共享侧栏母版组成，详见
表 \ref{tab:fe-pages}。

\begin{table}[h]
\centering
\caption{前端页面与功能}
\label{tab:fe-pages}
\small
\begin{tabularx}{\textwidth}{l l X}
\toprule
\textbf{页面} & \textbf{URL} & \textbf{功能} \\
\midrule
决策驾驶舱
  & \code{/}
  & 展示综合评分、4 项核心指标、今日决策摘要、3 套方案评分对比、
AI 解释区（"系统为什么这样排"） \\
智能排课
  & \code{/schedule.html}
  & 08:00--17:30 全日周课表（5 天 $\times$ 8 时段），
午休 12:00--13:30 显示为禁排灰带 \\
调课中心
  & \code{/reschedule.html}
  & 输入教师与时段，30 秒内返回最小影响新方案与变更明细 \\
资源概览
  & \code{/resources.html}
  & 教师 / 教室 / 班级 / 课程四张基础数据表 \\
\bottomrule
\end{tabularx}
\end{table}

\subsubsection{后端层}

后端由 6 个模块组成，模块划分与依赖关系如图
\ref{fig:arch} 所示。

\begin{itemize}[leftmargin=2em,itemsep=0.1em]
\item \code{\_\_init\_\_.py}：包标识；
\item \code{models.py}：不可变数据类（TimeSlot / Teacher / Room /
ClassGroup / CourseRequest / Problem / ScheduledLesson / SolveResult）；
\item \code{io.py}：JSON 与 \code{Problem} 之间的双向转换；
\item \code{validator.py}：业务数据与生成课表的双向校验；
\item \code{solver.py}：核心 CP-SAT 求解器（5 套策略）；
\item \code{service.py}：业务编排层（评分卡、调课、负载统计）；
\item \code{api.py}：HTTP 网关（6 个端点 + 静态文件服务）；
\item \code{cli.py}：命令行入口（用于 CI 验证）。
\end{itemize}

\subsection{数据流}

\begin{figure}[h]
\centering
\begin{tikzpicture}[
  node distance=0.8cm and 0.6cm,
  every node/.style={font=\small}
]
\node[cmblock] (json) {demo.json};
\node[cmblock, right=of json] (prob) {Problem};
\node[cmblock, right=of prob] (sol) {SolveResult};
\node[cmblock, right=of sol] (resp) {JSON Response};
\node[cmblock, below=of resp] (ui) {浏览器渲染};

\draw[cmarrow] (json) -- node[above]{load\_problem} (prob);
\draw[cmarrow] (prob) -- node[above]{solve\_schedule} (sol);
\draw[cmarrow] (sol) -- node[above]{to\_dict} (resp);
\draw[cmarrow] (resp) -- (ui);
\end{tikzpicture}
\caption{核心数据流}
\label{fig:dataflow}
\end{figure}

\subsection{模块接口设计}

\subsubsection{\code{models.Problem} 关键字段}

\begin{table}[h]
\centering
\caption{Problem 数据结构}
\label{tab:problem-struct}
\small
\begin{tabularx}{\textwidth}{l l l X}
\toprule
\textbf{字段} & \textbf{类型} & \textbf{示例} & \textbf{说明} \\
\midrule
\code{time\_slots} & \code{tuple[TimeSlot]} & 20 个
  & 含 \code{day / period / order} \\
\code{teachers} & \code{tuple[Teacher]} & 4 个
  & 学科资质 + 不可用 + 偏好 \\
\code{rooms} & \code{tuple[Room]} & 3 个
  & 容量 + 设备 + 不可用 \\
\code{classes} & \code{tuple[ClassGroup]} & 4 个
  & 人数 + 不可用 + 偏好 \\
\code{courses} & \code{tuple[CourseRequest]} & 4 个
  & 课次 + 合格教师 + 设备 + 允许时段 \\
\bottomrule
\end{tabularx}
\end{table}

\subsubsection{\code{SolveResult} 关键字段}

\begin{table}[h]
\centering
\caption{SolveResult 数据结构}
\label{tab:solve-result-struct}
\small
\begin{tabularx}{\textwidth}{l l X}
\toprule
\textbf{字段} & \textbf{类型} & \textbf{说明} \\
\midrule
\code{status}
  & \code{str}
  & \code{OPTIMAL} / \code{FEASIBLE} / \code{INFEASIBLE} / \code{INVALID\_DATA} \\
\code{schedule}
  & \code{list[ScheduledLesson]}
  & 求解得到的课次列表，按 slot.order 升序 \\
\code{conflicts}
  & \code{list[str]}
  & 独立校验器发现的冲突（应为空数组） \\
\code{metrics}
  & \code{dict}
  & 求解耗时、课次数、硬冲突数、目标值等 \\
\bottomrule
\end{tabularx}
\end{table}

\subsection{技术选型对比}

\begin{table}[h]
\centering
\caption{技术选型对比}
\label{tab:tech-choice}
\small
\begin{tabularx}{\textwidth}{l l l X}
\toprule
\textbf{决策点} & \textbf{选用方案} & \textbf{候选方案} & \textbf{选择理由} \\
\midrule
求解器
  & OR-Tools CP-SAT 9.15
  & Gurobi / 遗传算法 / 贪心
  & 免费、SOTA、硬约束严格 \\
后端框架
  & \code{BaseHTTPRequestHandler}
  & Flask / FastAPI / Django
  & 零依赖、跨平台、启动快 \\
前端框架
  & 原生 HTML/CSS/JS
  & Vue / React / Svelte
  & 零构建、零安装 \\
数据类
  & \code{dataclass(frozen=True)}
  & Pydantic / attrs
  & 标准库、轻量 \\
测试框架
  & \code{unittest}
  & pytest
  & 标准库、CI 无需额外安装 \\
CI
  & GitHub Actions
  & Travis CI / CircleCI
  & 免费、双版本、生态好 \\
\bottomrule
\end{tabularx}
\end{table}

\subsection{本章小结}

本章给出了系统的整体架构、模块划分、数据流与关键数据结构。
通过"前端 ↔ API ↔ Service ↔ Solver ↔ Validator"的分层设计，
保证了算法的可独立测试性与前端的可替换性。
后续章节将依次深入介绍算法设计（第 5 章）与系统实现
（第 6 章）。
